% GENERATED FILE. Edit canonical lessons, then run npm run generate:books.
% canonical-source-hash: fnv1a64:587d9457cc2314e1
% canonical-lessons: JA-C475-reizouko, JA-C475-shokki, JA-C475-ousaji, JA-C475-kozara, JA-C475-ajitsuke

\chapter{Kitchen and Table}
\label{ch:ja-kitchen-and-table}
\hlchaptermodality{475}

\begin{chapteropening}
\textbf{By the end of this chapter:} \emph{I can talk about the kitchen and the menu.}

Complete the last lesson of chapter 475: I can talk about the kitchen and the menu.
\end{chapteropening}

\section[\texorpdfstring{\ja{れいぞうこ} (reizōko)}{れいぞうこ (reizōko)}]{\ja{れいぞうこ} (reizōko) — a fridge; a refrigerator}

\label{lesson:JA-C475-reizouko}



\begin{quote}
Before the new one: say the Japanese for to tear into pieces, then the Japanese for to dilute.
\end{quote}

\subsection*{You'll want to know: \ja{れいぞうこ}}
\textbf{\ja{れいぞうこ}} — \emph{reizōko} — \textquotedblleft{}a fridge; a refrigerator\textquotedblright{}.

\begin{quote}
\textbf{\ja{ぎゅうにゅうは} \ja{れいぞうこに} \ja{あります。}} — \emph{gyūnyū wa reizōko ni arimasu} — \textquotedblleft{}The milk is in the fridge.\textquotedblright{}
\end{quote}

\begin{cousinweb}
\textbf{\ja{れいぞうこ}} (\emph{reizōko}) is Sino-Japanese (kango): three Chinese characters read with their on'yomi, readings borrowed from Chinese. They are \emph{rei}, \textquotedblleft{}cool\textquotedblright{} (Middle Chinese \emph{lengX}), \emph{zō}, \textquotedblleft{}storehouse\textquotedblright{} (\emph{dzangH}), and \emph{ko}, \textquotedblleft{}warehouse\textquotedblright{} (\emph{khuH}) (KANJIDIC2; Middle Chinese in Baxter's transcription).
\end{cousinweb}

\begin{grammarlens}[title={What you've built}]
One more word for describing what you lived and felt.
\end{grammarlens}

\subsection*{Your turn}
\begin{itemize}
\raggedright
  \item \emph{Say it:} \emph{reizōko}
  \item \emph{Say it:} \emph{reizōko}, once more
  \item \emph{Say it:} \emph{gyūnyū wa reizōko ni arimasu} — The milk is in the fridge.
  \item \emph{Recall:} say the Japanese for to tear into pieces, then the Japanese for to dilute, then say \emph{reizōko} again
\end{itemize}

\begin{tcolorbox}[breakable,colback=teal!4,colframe=teal!35!black,title={Before you move on}]
What does \ja{れいぞうこ} mean? (\textquotedblleft{}A fridge; a refrigerator\textquotedblright{}.) Say it once more.
\end{tcolorbox}
\section[\texorpdfstring{\ja{しょっき} (shokki)}{しょっき (shokki)}]{\ja{しょっき} (shokki) — tableware; the dishes}

\label{lesson:JA-C475-shokki}



\begin{quote}
Before the new one: say the Japanese for to dilute, then the Japanese for a fridge.
\end{quote}

\subsection*{You'll want to know: \ja{しょっき}}
\textbf{\ja{しょっき}} — \emph{shokki} — \textquotedblleft{}tableware; the dishes\textquotedblright{}.

\begin{quote}
\textbf{\ja{しょくじの} \ja{あとで、しょっきを} \ja{あらいました。}} — \emph{shokuji no ato de, shokki o araimashita} — \textquotedblleft{}After the meal I washed the dishes.\textquotedblright{}
\end{quote}

\begin{cousinweb}
\textbf{\ja{しょっき}} (\emph{shokki}) is Sino-Japanese (kango): two Chinese characters read with their on'yomi, readings borrowed from Chinese. They are \emph{sho}, \textquotedblleft{}eat\textquotedblright{} (Middle Chinese \emph{zyik}) and \emph{ki}, \textquotedblleft{}utensil\textquotedblright{} (\emph{khijH}) (KANJIDIC2; Middle Chinese in Baxter's transcription).
\end{cousinweb}

\begin{grammarlens}[title={What you've built}]
One more word for describing what you lived and felt.
\end{grammarlens}

\subsection*{Your turn}
\begin{itemize}
\raggedright
  \item \emph{Say it:} \emph{shokki}
  \item \emph{Say it:} \emph{shokki}, once more
  \item \emph{Say it:} \emph{shokuji no ato de, shokki o araimashita} — After the meal I washed the dishes.
  \item \emph{Recall:} say the Japanese for to dilute, then the Japanese for a fridge, then say \emph{shokki} again
\end{itemize}

\begin{tcolorbox}[breakable,colback=teal!4,colframe=teal!35!black,title={Before you move on}]
What does \ja{しょっき} mean? (\textquotedblleft{}Tableware; the dishes\textquotedblright{}.) Say it once more.
\end{tcolorbox}
\section[\texorpdfstring{\ja{おおさじ} (ōsaji)}{おおさじ (ōsaji)}]{\ja{おおさじ} (ōsaji) — a tablespoon}

\label{lesson:JA-C475-ousaji}



\begin{quote}
Before the new one: say the Japanese for a fridge, then the Japanese for tableware.
\end{quote}

\subsection*{You'll want to know: \ja{おおさじ}}
\textbf{\ja{おおさじ}} — \emph{ōsaji} — \textquotedblleft{}a tablespoon\textquotedblright{}.

\begin{quote}
\textbf{\ja{おおさじで} \ja{さとうを} \ja{いれて} \ja{ください。}} — \emph{ōsaji de satō o irete kudasai} — \textquotedblleft{}Please put the sugar in with a tablespoon.\textquotedblright{}
\end{quote}

\begin{cousinweb}
\textbf{\ja{おおさじ}} (\emph{ōsaji}) is a native Japanese compound (wago), inherited rather than borrowed from Chinese: \emph{ō}, \textquotedblleft{}big\textquotedblright{} and \emph{saji}, \textquotedblleft{}spoon\textquotedblright{} (JMdict, KANJIDIC2).
\end{cousinweb}

\begin{grammarlens}[title={What you've built}]
One more word for describing what you lived and felt.
\end{grammarlens}

\subsection*{Your turn}
\begin{itemize}
\raggedright
  \item \emph{Say it:} \emph{ōsaji}
  \item \emph{Say it:} \emph{ōsaji}, once more
  \item \emph{Say it:} \emph{ōsaji de satō o irete kudasai} — Please put the sugar in with a tablespoon.
  \item \emph{Recall:} say the Japanese for a fridge, then the Japanese for tableware, then say \emph{ōsaji} again
\end{itemize}

\begin{tcolorbox}[breakable,colback=teal!4,colframe=teal!35!black,title={Before you move on}]
What does \ja{おおさじ} mean? (\textquotedblleft{}A tablespoon\textquotedblright{}.) Say it once more.
\end{tcolorbox}
\section[\texorpdfstring{\ja{こざら} (kozara)}{こざら (kozara)}]{\ja{こざら} (kozara) — a small plate}

\label{lesson:JA-C475-kozara}



\begin{quote}
Before the new one: say the Japanese for tableware, then the Japanese for a tablespoon.
\end{quote}

\subsection*{You'll want to know: \ja{こざら}}
\textbf{\ja{こざら}} — \emph{kozara} — \textquotedblleft{}a small plate\textquotedblright{}.

\begin{quote}
\textbf{\ja{こざらに} \ja{しょうゆを} \ja{いれました。}} — \emph{kozara ni shōyu o iremashita} — \textquotedblleft{}I put soy sauce on a small plate.\textquotedblright{}
\end{quote}

\begin{cousinweb}
\textbf{\ja{こざら}} (\emph{kozara}) is a native Japanese compound (wago), inherited rather than borrowed from Chinese: \emph{ko}, \textquotedblleft{}small\textquotedblright{} and \emph{zara}, \textquotedblleft{}plate (sara)\textquotedblright{} (JMdict, KANJIDIC2).
\end{cousinweb}

\begin{grammarlens}[title={What you've built}]
One more word for describing what you lived and felt.
\end{grammarlens}

\subsection*{Your turn}
\begin{itemize}
\raggedright
  \item \emph{Say it:} \emph{kozara}
  \item \emph{Say it:} \emph{kozara}, once more
  \item \emph{Say it:} \emph{kozara ni shōyu o iremashita} — I put soy sauce on a small plate.
  \item \emph{Recall:} say the Japanese for tableware, then the Japanese for a tablespoon, then say \emph{kozara} again
\end{itemize}

\begin{tcolorbox}[breakable,colback=teal!4,colframe=teal!35!black,title={Before you move on}]
What does \ja{こざら} mean? (\textquotedblleft{}A small plate\textquotedblright{}.) Say it once more.
\end{tcolorbox}
\section[\texorpdfstring{\ja{あじつけ} (ajitsuke)}{あじつけ (ajitsuke)}]{\ja{あじつけ} (ajitsuke) — flavouring (a dish)}

\label{lesson:JA-C475-ajitsuke}



\begin{quote}
Before the new one: say the Japanese for a tablespoon, then the Japanese for a small plate.
\end{quote}

\subsection*{You'll want to know: \ja{あじつけ}}
\textbf{\ja{あじつけ}} — \emph{ajitsuke} — \textquotedblleft{}flavouring (a dish)\textquotedblright{}.

\begin{quote}
\textbf{\ja{ははの} \ja{あじつけは} \ja{うすいです。}} — \emph{haha no ajitsuke wa usui desu} — \textquotedblleft{}My mother flavours food lightly.\textquotedblright{}
\end{quote}

\begin{cousinweb}
\textbf{\ja{あじつけ}} (\emph{ajitsuke}) is a native Japanese compound (wago), inherited rather than borrowed from Chinese: \emph{aji}, \textquotedblleft{}taste\textquotedblright{} and \emph{tsuke}, \textquotedblleft{}applying (tsukeru)\textquotedblright{} (JMdict, KANJIDIC2).
\end{cousinweb}

\begin{grammarlens}[title={What you've built}]
One more word for describing what you lived and felt.
\end{grammarlens}

\subsection*{Your turn}
\begin{itemize}
\raggedright
  \item \emph{Say it:} \emph{ajitsuke}
  \item \emph{Say it:} \emph{ajitsuke}, once more
  \item \emph{Say it:} \emph{haha no ajitsuke wa usui desu} — My mother flavours food lightly.
  \item \emph{Recall:} say the Japanese for a tablespoon, then the Japanese for a small plate, then say \emph{ajitsuke} again
\end{itemize}

\begin{tcolorbox}[breakable,colback=teal!4,colframe=teal!35!black,title={Before you move on}]
What does \ja{あじつけ} mean? (\textquotedblleft{}Flavouring (a dish)\textquotedblright{}.) Say it once more.
\end{tcolorbox}
